\documentclass[11pt]{article}
\usepackage[T1]{fontenc}
\usepackage{lmodern}
\usepackage[margin=1in]{geometry}
\usepackage{amsmath,amssymb,amsthm,booktabs}
\usepackage[hidelinks]{hyperref}
\clubpenalty=10000
\widowpenalty=10000
\newtheorem{theorem}{Theorem}[section]
\newtheorem{proposition}[theorem]{Proposition}
\newtheorem{corollary}[theorem]{Corollary}
\newtheorem{lemma}[theorem]{Lemma}
\theoremstyle{definition}
\newtheorem{definition}[theorem]{Definition}
\newtheorem{remark}[theorem]{Remark}
\newcommand{\NN}{\mathbb N}
\newcommand{\QQ}{\mathbb Q}
\newcommand{\vp}{v_p}
\title{Effective Denominator Primes in Rational Semirings:\\
Exact Criteria, Constructive Certificates, and a Unit Obstruction}
\author{Rational Sumring Lab\\\small Working research note; author attribution to be reviewed}
\date{November 1, 2026}
\begin{document}
\maketitle
\begin{abstract}
For positive rational numbers $q,r$, we study the primes that occur in the
reduced denominator of a monomial $q^ir^j<1$. We give an elementary exact
criterion valid for every positive rational pair, including pairs with both
generators below one. In the mixed case, we prove the criterion without
logarithmic approximation and give an explicit rational bound for constructing
a witness. We show that the same prime set is detected by all semiring elements
below one. A separate valuation obstruction proves that the three examples
highlighted in CrowdMath 2026 Open Problem 2 have trivial multiplicative unit
groups, and we identify a multiplicative atom in each of those examples.
These results support exact software certificates. They do not decide
the atomicity or irreducible-divisor-finiteness questions in that problem.
No priority or novelty claim is made for the elementary results proved here.
\end{abstract}

\section{Scope, definitions, and provenance}
For $q,r\in\QQ_{>0}$, put
\[
 S(q,r)=\NN_0[q,r]
 =\left\{\sum_{i,j\ge0}c_{ij}q^ir^j:
       c_{ij}\in\NN_0,\ \text{finitely many }c_{ij}\ne0\right\},
 \qquad M(q,r)=S(q,r)\setminus\{0\}.
\]
Here $M(q,r)$ is a commutative cancellative monoid under multiplication.
Its units are the elements $u$ such that $u^{-1}\in M(q,r)$. Elements are
associates if they differ by a unit. An atom is a nonunit $a$ for which
$a=bc$ in $M(q,r)$ implies that $b$ or $c$ is a unit. A monoid is atomic if
each nonunit is a finite product of atoms; it is irreducible-divisor finite
(IDF) if each element has only finitely many nonassociate atom divisors.
All statements about atoms and IDF in this note concern multiplication,
not addition.

The motivation is the official \emph{Arithmetic Conjectures on Rational
Semirings} project, especially Open Problem 2~\cite{cm2}.
Its effective-prime definition is stated for $0<q<1<r$; we explicitly extend
that definition to all positive pairs. The repository retains its historical
name ``Rational Sumring Lab.''

\begin{definition}
For a prime $p$, let $v_p$ be the integer-valued valuation on $\QQ_{>0}$,
normalized by $v_p(p)=1$. Define
\[
 E(q,r)=\{p\text{ prime}:\text{ some }i,j\in\NN_0
       \text{ satisfy }q^ir^j<1\text{ and }v_p(q^ir^j)<0\}.
\]
For a positive rational $x$, write $D(x)$ for the prime divisors of its
denominator in lowest terms.
\end{definition}

The proofs below are self-contained elementary derivations. The mixed-case
criterion already appeared in this repository's earlier proof note; it is
not presented here as a newly discovered theorem. The additional proofs,
constructive bound, and implementation repair have not been shown to be new
relative to the literature or the CrowdMath discussion. This is a
verification-oriented working note, not a claim to have solved an open problem.
The draft and accompanying software were prepared with AI assistance and
require the named authors' mathematical and attribution review before any
submission. No affiliation with, or endorsement by, MIT or CrowdMath is claimed.

\section{An exact criterion and explicit witnesses}
The elementary identities used throughout are
\[
 v_p(xy)=v_p(x)+v_p(y),\qquad
 v_p(x_1+\cdots+x_n)\ge\min_k v_p(x_k).
\]
The second follows by factoring out the smallest common power of $p$ and
putting the remaining summands over denominators prime to $p$.

\begin{theorem}[Mixed-case criterion]\label{thm:mixed}
Suppose $0<q<1<r$, and fix a prime $p$. Set $A=v_p(q)$ and $C=v_p(r)$.
Then
\[
 p\in E(q,r)\quad\Longleftrightarrow\quad
 A<0\quad\text{or}\quad
 \bigl(A\ge0,\ C<0,\ q^{-C}r^A<1\bigr).
\]
\end{theorem}
\begin{proof}
If $A<0$, the monomial $q$ is a witness. If $A,C\ge0$, every monomial has
nonnegative valuation, so there is no witness.

It remains to treat $A\ge0$ and $C=-B<0$, where $B$ is a positive integer.
Put $t=q^Br^A$. For any witness $m=q^ir^j$, we have $i>0$ and $Bj-Ai>0$.
If $t\ge1$, then
\[
 m^B=q^{Bi}r^{Bj}=t^i r^{Bj-Ai}>1,
\]
contrary to $m<1$. Thus $t<1$ is necessary. If $t<1$, choose a positive
integer $k$ with $rt^k<1$. Such a $k$ exists since $t^k\to0$.
Taking $i=Bk$ and $j=Ak+1$ gives
\[
 q^ir^j=rt^k<1,\qquad Ai-Bj=-B<0.
\]
The next proposition gives an explicit choice of $k$ using rational
arithmetic alone.
\end{proof}

\begin{proposition}[A rational witness bound]\label{prop:bound}
In the final case of Theorem~\ref{thm:mixed}, suppose $t=q^Br^A<1$. Set
\[
 K=\left\lfloor\frac{(r-1)t}{1-t}\right\rfloor+1.
\]
Then $K\ge1$, and $(i,j)=(BK,AK+1)$ is a witness. In particular, the least
positive integer $k$ with $rt^k<1$ satisfies $k\le K$.
\end{proposition}
\begin{proof}
Let $\delta=t^{-1}-1=(1-t)/t>0$. The definition of $K$ gives
$K\delta>r-1$. The binomial theorem implies
\[
 t^{-K}=(1+\delta)^K\ge1+K\delta>r,
\]
so $rt^K<1$. Its valuation at the displayed exponent pair is $-B$, as in
the preceding proof. The strict inequality in the definition of $K$ is
essential; replacing it by a potentially equal endpoint is not justified.
\end{proof}

\begin{corollary}[All positive rational pairs]\label{cor:all}
The set $E(q,r)$ is determined as follows.
\begin{enumerate}
\item If $q,r\ge1$, then $E(q,r)=\varnothing$.
\item If $q,r\le1$, then $E(q,r)=D(q)\cup D(r)$.
\item If one generator is below one and the other is above one, reorder
them as $u<1<v$ and apply Theorem~\ref{thm:mixed} to $(u,v)$.
\end{enumerate}
The first two cases agree at $(q,r)=(1,1)$. Cases with one generator equal
to one are included. Only primes in $D(q)\cup D(r)$ need to be tested.
\end{corollary}
\begin{proof}
In the first case no monomial is below one. In the second case any
denominator prime of either generator is witnessed by that generator:
a generator with a nontrivial denominator is strictly below one in this
case. Conversely, outside $D(q)\cup D(r)$ both generator valuations are
nonnegative. This also proves the final assertion. Swapping generators
swaps the witness exponents and changes neither the semiring nor $E$.
\end{proof}

\begin{remark}[What the bound does and does not provide]
Proposition~\ref{prop:bound} proves termination, not a polynomial bit-complexity
bound. The integer $K$ can be large when $t$ is close to one, and the
numerator and denominator of a certificate can be large. One may find the
least $k$ in the special family $(Bk,Ak+1)$ by doubling an upper bound and
then binary search on the monotone predicate $rt^k<1$. This does not find
a globally smallest exponent pair. All comparisons remain exact.
\end{remark}

\section{Denominators of sums and an obstruction to units}
\begin{proposition}\label{prop:sums}
A prime $p$ lies in $E(q,r)$ if and only if some $x\in S(q,r)$ satisfies
$0<x<1$ and $v_p(x)<0$. If $p\in E(q,r)$, then for every $\varepsilon>0$
and positive integer $L$, some $x\in S(q,r)$ satisfies
$0<x<\varepsilon$ and $v_p(x)\le-L$.
\end{proposition}
\begin{proof}
A monomial witness is itself such an element. Conversely, write $x$ as a
finite sum of nonzero terms $c_{ij}q^ir^j$, with $c_{ij}\in\NN$.
The valuation inequality forces at least one term to have negative
valuation. Since $v_p(c_{ij})\ge0$, its monomial has negative valuation.
Positivity and $c_{ij}\ge1$ give $q^ir^j\le c_{ij}q^ir^j\le x<1$.
This monomial witnesses effectiveness. Finally, powers $m^n$ of any
monomial witness tend to zero, and $v_p(m^n)=nv_p(m)\to-\infty$.
\end{proof}

\begin{proposition}[Sufficient condition for trivial units]\label{prop:units}
Suppose a prime $\ell$ satisfies
\[
 v_\ell(q)\ge0,\quad v_\ell(r)\ge0,\qquad
 v_\ell(g)>0\text{ for every }g\in\{q,r\}\text{ with }g<1.
\]
Then $M(q,r)^\times=\{1\}$.
\end{proposition}
\begin{proof}
Every nonzero element of $S(q,r)$ has nonnegative $\ell$-valuation.
A monomial below one must involve a generator below one with positive
exponent. The hypotheses therefore give valuation at least one for
each such monomial. If $0<x<1$ belongs to $S(q,r)$, every nonzero term
in a polynomial representation is below one. As in
Proposition~\ref{prop:sums}, its underlying monomial is below one.
Every term has valuation at least one, so $v_\ell(x)\ge1$.
If a unit $u\ne1$ existed, either $u$ or $u^{-1}$ would be below one.
Its inverse would have negative $\ell$-valuation, contradicting membership
in $S(q,r)$. Thus the only unit is $1$.
\end{proof}

We use this as a sufficient criterion and make no classification of unit
groups for arbitrary pairs.

\section{The three specified examples}
The following table records results proved here. The auxiliary prime
$\ell$ certifies the unit group by Proposition~\ref{prop:units}; it is
not required to be a denominator prime.
\begin{center}
\begin{tabular}{@{}llll@{}}
\toprule
$(q,r)$ & $E(q,r)$ & $\ell$ & $M(q,r)^\times$\\
\midrule
$(6/35,19/3)$ & $\{5,7\}$ & $2$ & $\{1\}$\\
$(3/10,7/5)$ & $\{2,5\}$ & $3$ & $\{1\}$\\
$(2/5,2/3)$ & $\{3,5\}$ & $2$ & $\{1\}$\\
\bottomrule
\end{tabular}
\end{center}

For the first pair, $q$ witnesses $5$ and $7$. The only remaining
denominator prime is $3$, with $A=1$, $B=1$, and
$t=qr=114/105=38/35>1$, so it is ineffective. Also
$v_2(q)=1$ and $v_2(r)=0$, which verifies the unit obstruction.
For the second pair, $q$ witnesses both denominator primes $2,5$;
$v_3(q)=1$ and $v_3(r)=0$ verify the unit obstruction.
For the third pair, both generators are below one: $r$ witnesses $3$
and $q$ witnesses $5$, while both have $2$-valuation one.
Thus associates are equal in all three monoids. Nothing in these arguments
establishes that their nonunits factor into atoms.

\begin{proposition}[A specified atom in each example]\label{prop:atom}
In each row of the table, the generator $q$ is a multiplicative atom.
In $M(2/5,2/3)$ the other generator is reducible:
\[
 \frac23=\frac25\left(1+\frac23\right).
\]
\end{proposition}
\begin{proof}
Use the prime $\ell$ in the corresponding row. In every case
$v_\ell(q)=1$, every element has nonnegative valuation, and every element
below one has strictly positive valuation, as proved above.
We first show that an element $s<1$ of valuation one must satisfy $s\ge q$.
For either mixed pair, a representation of $s$ can contain only monomials
with $q$-exponent $i\ge1$, since $r>1$. If every term had $i\ge2$, then
each term and hence the sum would have valuation at least two. Thus some
term has $i=1$; since $r\ge1$ and its coefficient is positive, that term
is at least $q$. For the third pair, a constant term is impossible and
the valuation of a monomial is its total degree. A sum of valuation one
must therefore contain a degree-one term, which is at least
$\min(q,r)=q$. This proves the claim in all cases.

If $q=ab$ in $M(q,r)$, the nonnegative integers
$v_\ell(a),v_\ell(b)$ sum to one. Name the zero-valuation factor $a$.
It cannot be below one, so $a\ge1$. If $a>1$, then $b=q/a<q<1$ and
$v_\ell(b)=1$, contradicting the preceding claim. Therefore $a=1$.
The element $q$ is a nonunit by Proposition~\ref{prop:units}, so it is an atom.
Finally, the displayed factorization of $2/3$ has both factors in the
semiring and neither equal to one. Both are nonunits by the unit result.
\end{proof}

Identifying one atom is not an atomicity theorem. The proposition leaves
the behavior of all other elements and all atom divisors unclassified.

As a nontrivial construction example, take $(q,r,p)=(4/5,11/8,2)$.
Then $A=2$, $B=3$, $t=121/125$, and Proposition~\ref{prop:bound} gives
$K=12$. Hence
\[
 (i,j)=(36,25),\qquad
 q^{36}r^{25}=\frac{11^{25}}{8\cdot5^{36}}<1,\qquad
 v_2(q^{36}r^{25})=-3.
\]
The rational bound need not be sharp. For $(q,r)=(2/3,3/2)$ the
candidate $p=2$ has $A=B=1$ and $t=1$, so it is ineffective;
$p=3$ is effective via $q$. This boundary example checks the strict
inequality in Theorem~\ref{thm:mixed}.

\section{Relation to the open questions}
Open Problem 1~\cite{cm1} asks for a classification of rational pairs with
atomic non-IDF or IDF non-atomic multiplicative monoids. No such
classification is supplied here; our statements concern positive pairs.

For comparison, Deng, Gotti, and Zeng~\cite[Theorem 5.7]{dgz} classify
the \emph{multiplicative finite factorization} property for a single
positive rational generator $q$: it holds exactly when $q\ge1$, $q$
is a unit fraction, or its reduced denominator is a prime power.
This established one-generator result does not give a classification
for two generators. Their Theorem 5.8 relates bounded factorizations
and ACCP. Neither failure of finite factorization nor failure of ACCP
alone means failure of atomicity.

Dani, Gotti, Li, and Paladiya~\cite[Introduction, footnote 1]{dglp}
report forthcoming work of Cen et al. on non-atomic rational monogenic
examples. The cited footnote supplies no proof, and a complete public
proof was not located in this audit. We record it as an announcement,
not a theorem verified here. It also does not answer the different
existence-of-atomic-examples question in Open Problem 2(c).

The project statement~\cite{cm2} reports prior message-board work:
for a multiplicatively independent mixed pair, at least two effective
primes together with an ineffective prime divisor of the numerator of $q$
imply failure of the ascending chain condition on principal ideals (ACCP),
and failure of atomicity or IDF. That prior result is attributed to the
project; it is not reproved here or counted as a contribution of this note.
Both mixed examples satisfy its stated hypotheses. Their independence
can be checked by valuations: $v_2$ and $v_{19}$ separate the generators
of the first pair, and $v_3$ and $v_7$ separate those of the second.

The questions recorded in Open Problem 2 ask which property fails for
the two mixed examples, seek a broader classification, ask whether some
monogenic example with numerator greater than one and at least two
denominator primes is atomic, and ask about atomicity and IDF for
$\NN_0[2/5,2/3]\setminus\{0\}$. This note answers none of these questions.
The disjunction ``not atomic or not IDF'' does not decide either property
separately. A finite collection of divisor witnesses cannot decide them.

An atomicity proof would need a factorization argument for every nonunit;
a non-atomicity proof would need an element with no finite atom
factorization. An IDF proof must control all nonassociate atom divisors of
each element. To disprove IDF, one must exhibit an element with infinitely
many such atom divisors and prove their irreducibility and
nonassociateness. Trivial units simplify the last issue but settle none
of the other requirements. These are the missing proof obligations.

\section{Reproducible implementation and verification}
The companion code uses integer arithmetic and Python's
\texttt{fractions.Fraction}. The exact prime decision tests only the
denominator-prime union and handles all cases of Corollary~\ref{cor:all}.
A positive certificate records a prime, nonnegative integer exponents,
and the exact monomial value. Verification checks primality, the input
domain, the exponent domain, equality with the recomputed monomial, and
both strict inequalities. Acceptance proves effectiveness for that
prime; completeness of the decision procedure follows from the theorem,
not from the verifier alone.

The repository's original failure on $(2/5,2/3)$ was a domain mismatch:
tests included this pair but the exact and bounded routines required
$0<q<1<r$. Corollary~\ref{cor:all} supplies the missing case. The
corrected tests cover boundary equality, swapped inputs, generators equal
to one, invalid inputs, malformed certificates, and independent bounded
witness checks. The separate script \texttt{paper/verify\_results.py}
checks the displayed examples and a finite rational grid for the explicit
bound. These are executable consistency checks, not formal proof-assistant
verification or a substitute for the proofs above.

Finite enumeration is explicitly budgeted. It certifies membership of
the elements it produces and factorizations whose factors are present.
Failure to find a factorization leaves an unverified candidate, not an
atom. Moreover, a factor unequal to one need not be a nonunit in a general
semiring; unit information must be supplied or proved separately.
No empirical search in this software is interpreted as a resolution of
atomicity or IDF.

\begin{thebibliography}{9}
\bibitem{dgz}
Anna Deng, Felix Gotti, and Jason Zeng,
\emph{Factorizations in rational monogenic semidomains},
arXiv:2607.10178 (2026).
\url{https://arxiv.org/abs/2607.10178}.
\bibitem{dglp}
Jiya Dani, Felix Gotti, Bryan Li, and Arav Paladiya,
\emph{Factorizations in algebraic monogenic semidomains},
arXiv:2609.32138 (2026).
\url{https://arxiv.org/abs/2609.32138}.
\bibitem{cm2}
MIT PRIMES / Art of Problem Solving CrowdMath,
\emph{Arithmetic Conjectures on Rational Semirings, Open Problem 2},
updated October 15, 2026; consulted November 1, 2026.
\url{https://artofproblemsolving.com/polymath/mitprimes2026/p/open-prob2}.
\bibitem{cm1}
MIT PRIMES / Art of Problem Solving CrowdMath,
\emph{Arithmetic Conjectures on Rational Semirings, Open Problem 1},
consulted November 1, 2026.
\url{https://artofproblemsolving.com/polymath/mitprimes2026/p/prob1}.
\end{thebibliography}
\end{document}
